\clearpage
\SourceLeaf{264}
\begin{center}
{\large ABSCHNITT IV.}\par\medskip
\pdfbookmark[1]{Über die Bahnbestimmung unter Berücksichtigung von Störungen.}{liber-ii-sectio-iv}
\textit{Über die Bahnbestimmung unter Berücksichtigung von Störungen.}
\end{center}
\Needspace{6\baselineskip}
\begin{center}\hypertarget{native-264-article190}{}\pdfbookmark[2]{190.}{article-190}190.\end{center}
Die Störungen, die die Bewegungen der Planeten durch die Wirkung der
übrigen Planeten erfahren, sind so gering und vollziehen sich so langsam,
dass sie erst nach längerer Zeit merklich werden. Über einen kürzeren
Zeitraum --- oder je nach den Umständen sogar über einen oder mehrere
volle Umläufe --- weicht die Bewegung so wenig von der Bewegung auf
einer vollkommenen, genau nach den Gesetzen \textsc{Kepler}s durchlaufenen
Ellipse ab, dass die Abweichung aus den Beobachtungen nicht hervorgeht.
Solange dies der Fall ist, lohnt es sich kaum, die Störungen vorzeitig
zu berechnen. Es genügt vielmehr, einen gleichsam oskulierenden Kegelschnitt
an die Beobachtungen anzupassen. Wenn der Planet längere Zeit genau
beobachtet worden ist, macht sich die Wirkung der Störungen schließlich
so deutlich bemerkbar, dass sich nicht mehr alle Beobachtungen mit einer
rein elliptischen Bewegung genau in Einklang bringen lassen. Eine
vollständige und dauerhafte Übereinstimmung lässt sich dann nur erzielen,
wenn man die Störungen mit der elliptischen Bewegung in geeigneter Weise
verbindet.

Die Bestimmung der elliptischen Elemente, mit denen die Störungen zu
verbinden sind, um die Beobachtungen genau darzustellen, setzt die
Kenntnis dieser Störungen voraus. Umgekehrt kann man die Theorie der
Störungen nur dann genau ausarbeiten, wenn die Elemente bereits
näherungsweise bekannt sind. Die Natur der Sache lässt es nicht zu, diese
schwierige Aufgabe gleich beim ersten Versuch vollkommen zu lösen. Die
höchste Genauigkeit bei der Bestimmung der Störungen und der Elemente
lässt sich erst durch mehrmalige abwechselnde Korrektur erreichen. Für
die erste Ausarbeitung der Störungstheorie legt man daher die rein
elliptischen \SourceLeaf{265}Elemente zugrunde, die näherungsweise an die
Beobachtungen angepasst worden sind.
Dann sucht man eine neue Bahn, die zusammen mit diesen Störungen den
Beobachtungen möglichst genau entspricht. Weicht sie erheblich von der
vorherigen Bahn ab, so sind die Störungen auf dieser Grundlage erneut
zu berechnen. Diese abwechselnden Korrekturen werden so oft wiederholt,
bis die Beobachtungen, die Elemente und die Störungen möglichst genau
übereinstimmen.
